\documentclass[12pt]{article}
\usepackage[utf8]{inputenc}
\usepackage[T1]{fontenc}
\usepackage{hyperref}
\title{Tobanga Language Resources for K-12 Curriculum Development}
\author{Compiled by Claude Sonnet 5 (claude-sonnet-5) for Andr\'as Kornai}
\date{}
\begin{document}
\maketitle

\section{Introduction}
Tobanga (ISO 639-3 tng; Glottocode toba1280), also called Northern Gabri, is an East Chadic (Afro-Asiatic) language of the Tandjil\'e Region of Chad. Sources cited in the search give about 30,000 speakers (1999 figure). Joshua Project describes the Tobanga as a Baggara subgroup, mostly herdsmen, and gives population figures that differ between listings. Few resources turned up in the searches, and none are educational.

\section{Primary Educational Resources}
No primers, textbooks or school programmes were found.

\section{Linguistic and Reference Materials}
\begin{itemize}
\item Wikipedia, Tobanga language: \url{https://en.wikipedia.org/wiki/Tobanga_language}
\item ISO 639 reference entry for tng: \url{https://wikipedia.com/wiki/ISO_639:tng}
\end{itemize}

\section{Religious and Community Resources}
\begin{itemize}
\item Joshua Project, Tobanga in Chad: \url{https://joshuaproject.net/people_groups/15520}
\item Tobanga New Testament (1978) on YouVersion: \url{https://www.bible.com/versions/3711}
\item Tobanga New Testament on Bibliatodo: \url{https://www.bibliatodo.com/ht/amp/bib-la/vesyon/tng1978}
\item ScriptureEarth page for tng: \url{https://www.scriptureearth.org/00eng.php?iso=tng}
\end{itemize}

\section{Digital and Audio Resources}
\begin{itemize}
\item Global Recordings Network page for tng (audio): \url{https://www.globalrecordings.net/en/language/tng}
\end{itemize}

\section{Curriculum Development Context}
The language is stable and widely used in its community, but the searches turned up no orthography or educational materials. The 1978 New Testament is the main written text and could be a starting corpus. A curriculum would have to be built from new documentation. It would have to be coordinated with Chad's language-in-education policy, which I did not research here.
\end{document}